\TieudeT{PHONG TỎA VẬT LÍ}
\subsubsection*{PHẦN I. Thí sinh trả lời câu hỏi từ câu 1 đến câu 18. Mỗi câu hỏi thí sinh chỉ chọn một phương án}
\setcounter{ex}{0}
%mac
%\loai{Ném xuống - tính thời gian chạm đất}
\begin{ex}%Tailieuchuan
Một học sinh đứng lan can tầng bốn ném quả cầu thẳng đứng lên trên, tiếp theo đó ném tiếp quả cầu thẳng đứng xuống dưới với cùng tốc độ. Bỏ qua sức cản của không khí, quả cầu nào chạm mặt đất có tốc độ lớn hơn?
\choice 
{Quả cầu ném lên}
{Quả cầu ném xuống}
{\True Cả hai quả cầu chạm đất có cùng tốc độ}
{Không xác định được vân tốc quả cầu vì thiếu độ cao}
\loigiai{}
\end{ex}
\begin{ex}
Một vật được ném thẳng đứng lên cao với tốc độ ban đầu $20\ m/s$. Bỏ qua sức cản không khí, lấy $g = 10\ m/s^2$. Thời gian vật lên tới độ cao cực đại là
\choice
{\True $2\ s$}
{$1\ s$}
{$4\ s$}
{$0,5\ s$}
\loigiai{}
\end{ex}

\begin{ex}
Một vật được ném thẳng đứng lên cao với vận tốc $30\ m/s$. Lấy $g = 10\ m/s^2$. Độ cao cực đại vật đạt được so với vị trí được ném là
\choice
{$60\ m$}
{\True $45\ m$}
{$30\ m$}
{$15\ m$}
\loigiai{}
\end{ex}

\begin{ex}
Một vật được ném thẳng đứng xuống dưới với vận tốc ban đầu $10\ m/s$ từ độ cao $80\ m$. Bỏ qua sức cản không khí, lấy $g = 10\ m/s^2$. Thời gian từ lúc ném vật đến khi vật chạm đất là
\choice
{$3,123\ s$}
{\True $4,462\ s$}
{$5,563\ s$}
{$6,634\ s$}
\loigiai{}
\end{ex}

\begin{ex}
Một vật được ném thẳng đứng lên cao từ mặt đất. Sau $2\ s$ thì vận tốc của vật là $0$. Lấy $g = 10\ m/s^2$. Vận tốc ban đầu của vật là
\choice
{$5\ m/s$}
{\True $20\ m/s$}
{$10\ m/s$}
{$15\ m/s$}
\loigiai{}
\end{ex}

\begin{ex}
Một vật rơi tự do từ độ cao $45\ m$ so với mặt đất. Bỏ qua sức cản không khí, lấy $g = 10\ m/s^2$. Thời gian rơi là
\choice
{$2\ s$}
{$4\ s$}
{\True $3\ s$}
{$4,5\ s$}
\loigiai{}
\end{ex}

\begin{ex}
Một vật được ném thẳng đứng lên cao với vận tốc ban đầu $v_0$. Bỏ qua sức cản không khí, thời gian vật lên đến độ cao cực đại là $t$. Thời gian rơi từ đó xuống lại vị trí ném là
\choice
{$\dfrac{t}{2}$}
{$2t$}
{\True $t$}
{$3t$}
\loigiai{}
\end{ex}

\begin{ex}
Một vật rơi tự do từ độ cao $20\ m$. Lấy $g = 10\ m/s^2$. Tốc độ ngay trước khi chạm đất là
\choice
{$10\ m/s$}
{$15\ m/s$}
{\True $20\ m/s$}
{$25\ m/s$}
\loigiai{}
\end{ex}

\begin{ex}
Một vật được ném thẳng đứng lên cao. Khi đến độ cao $h = 20\ m$ thì vận tốc của vật là $10\ m/s$. Lấy $g = 10\ m/s^2$. Tốc độ ban đầu của vật là
\choice
{\True $22,36\ m/s$}
{$30,34\ m/s$}
{$25\ m/s$}
{$15\ m/s$}
\loigiai{}
\end{ex}
\begin{ex}%Tailieuchuan
Một người đứng trên tòa nhà có độ cao $120\ m$, ném một vật thẳng đứng xuống dưới với vận tốc $10 \ m/s$. Cho $g=10 \ m/s^2$. Tốc độ của vật ngay trước khi chạm đất là
\choice
{$20 \ m/s$}
{$30 \ m/s$}
{$40 \ m/s$}
{\True $50\ m/s$}
\loigiai{}
\end{ex}
\begin{ex}%Tailieuchuan
Một người đứng trên tòa nhà có độ cao $120\ m$, ném một vật thẳng đứng xuống dưới với vận tốc $10 \ m/s$ cho $g=10 \ m/s^2$. Kể từ lúc ném, vật chạm đất sau			
\choice
{\True $4\ s$}
{$5\ s$}
{$6\ s$}
{$7\ s$}
\loigiai{}
\end{ex}
\begin{ex}%book %[1P1K7-4][Loai1] %BAI 4+5 - HOC MAI
Một hòn sỏi nhỏ được ném thẳng đứng xuống dưới với vận tốc đầu bằng $9,8\ m/s$ từ độ cao $39,2\ m$. Lấy $g = 9,8\  m/s^2$. Bỏ qua lực cản của không khí. Sau bao lâu hòn sỏi rơi tới đất?
\choice
{$t=1\ s$}
{\True $t=2\ s$}
{$t=3\ s$}
{$t=4\ s$}
\haidong{14}
\loigiai{
\begin{itemize}
\item Ta có gọi $v_t$ là vận tốc cuối khi chạm đất $v^2_t - v^2_0 = 2gh$ $\Rightarrow$ $v^2_t = 9,8^2 + 2.9,8.39,2 \Rightarrow v_t = 29,4$  m/s.
\item Thời gian hòn đá rơi xuống đất là: $t = \dfrac{v_t - v_0}{g} = \dfrac{29,4 - 9,8}{9,8} = 2$ s.
\end{itemize}	}
\end{ex}
%\loai{Ném lên - tính $h_{\max}$}
\begin{ex}%book %[1P1B7-4][Loai1] %BAI 4+5 - HOC MAI
Một vật ném thẳng đứng lên cao với tốc độ  ban đầu $36\ km/h$. Lấy $g=10\ m/s^2$. Độ cao cực đại $h_{\max}$ so với vị trí ném mà vật có thể đạt tới là
\choice
{$h_{\max} = 15\ m$}
{$h_{\max} = 10\ m$}
{\True $h_{\max} = 5\ m$}
{$h_{\max} = 0,5\ m$}
\haidong{14}
\loigiai{

\begin{itemize}
\item Đổi $36\  km/h = 10\  m/s$.
\item Áp dụng các công thức tính $v^2_t - v^2_0 = 2ah$ 
\item Khi vật lên cao nhất thì $v_t = 0 ;\ a = -g ;\ h = h_{\max} \Rightarrow h_{\max} = \dfrac{v^2_0}{2g} = 5$ m.
\end{itemize}	}
\end{ex}
%\loai{Tính vận tốc ném}
\begin{ex}%book %[1P1B7-2][Loai1] %BAI 4+5 - HOC MAI
Một người thợ xây ném một viên gạch theo phương thẳng đứng cho một người khác ở trên tầng cao 4 m. Người này chỉ việc giơ tay ngang ra là bắt được viên gạch. Lấy $g = 9,8\  m/s^ 2$. Tốc độ khi ném là bao nhiêu để cho vận tốc viên gạch lúc người kia bắt được là bằng $0$?
\choice
{\True $8,85\  m/s$}
{$8,94\  m/s$}
{$6,26\  m/s$}
{$6,32\  m/s$}
\haidong{14}
\loigiai{
\begin{itemize}
\item Khi vận tốc viên gạch bằng 0 thì viên gạch đạt độ cao cực đại 4 m. 
\item Vận tốc ban đầu của viên gạch được tính theo công thức: $v_0=\sqrt{2gh}=\sqrt{2.9,8.4}=8,85\  m/s$.
\end{itemize}	}
\end{ex}
%\loai{Ném lên - tính h khi cho v}
\begin{ex}%book %[1P1B7-4][Loai1] %BAI 4+5 - HOC MAI
Một người ném một quả bóng từ mặt đất lên cao theo phương thẳng đứng với tốc độ $4\  m/s$. Lấy $g = 9,8\  m/s^2$. Độ cao của quả bóng khi nó có tốc độ $2,5\ m/s$ là 
\choice
{\True $0,497\ m$}
{$0,994\ m$}
{$0,077\ m$}
{$0,153\ m$}
\haidong{14}
\loigiai{
Độ cao của quả bóng là: $h=s=\dfrac{v_t^2-{v_0}^2}{2a}=\dfrac{2,5^2-4^2}{2.(-9,8)}=0,497\ m$.
}
\end{ex}
%\loai{Thời gian từ lúc ném đến khi chạm đất}
\begin{ex}%book %[1P1K7-4][Loai1] %BAI 4+5 - HOC MAI
Người ta ném một vật từ mặt đất lên trên cao theo phương thẳng đứng với tốc độ $4\ m/s$. Lấy $g = 9,8\  m/s^2$. Sau bao lâu thì vật đó rơi chạm đất?
\choice
{$0,2\ s$}
{$0,41\ s$}
{\True $0,82\ s$}
{$1,64\ s$}
\haidong{14}
\loigiai{

\begin{itemize}
\item Chọn trục tọa độ có phương thẳng đứng và hướng lên trên, gốc O tại mặt đất. Phương trình chuyển động là: $y=4t - \dfrac{1}{2}.9,8t^2 $
\item Khi vật chạm đất, y = 0 ta có: $0=4t - \dfrac{1}{2}.9,8t^2 \Rightarrow t=0$ và $t=\dfrac{4}{{4,9}}\approx 0,82\ s$
\item Vậy sau 0,82 s vật rơi xuống đất.
\end{itemize}	}
\end{ex}
%\loai{Ném lên - tính vận tốc chạm đất}
\begin{ex}%book
Một vật được ném thẳng đứng lên trên. Một người quan sát thấy vật đi qua độ cao $40\ m$ hai lần cách nhau $2\ s$. Bỏ qua sức cản của không khí, lấy $g=10\ m/s^2$. Vận tốc ban đầu khi ném vật là
\choice
{$15\ m/s$}
{\True $30\ m/s$}
{$45\ m/s$}
{$60\ m/s$}

\haidong{17}
\loigiai{
Chọn trục tọa độ $Oy$ thẳng đứng hướng xuống gốc toạ độ O tại vị trí có độ cao $40\ m$.

$t_1$ là thời gian vật đi lên đến vị trí cao nhất.

$t_2$ là thời gian vật đi từ vị trí cao nhất đến O.

$\left\{\begin{array}{l}t_1=\sqrt{\dfrac{2 h}{g}} \\ t_1=\sqrt{\dfrac{2 h}{g}}\end{array} \rightarrow t_1 t_2=\dfrac{2 h}{g}\right.$ ta có $t_1 t_2=\dfrac{2 h}{g} \Rightarrow t_1 t_2=\dfrac{240}{10} \Rightarrow t_1 t_2=8$

Bài toán cho $t_2-t_1=2$

Vậy $t_1=2 s, t_2=4 s$

Mặt khác $t_1+t_2=\dfrac{2 v_0}{g} \Leftrightarrow t_1+t_2=\dfrac{2 v_0}{10} \Rightarrow v_0=5\left(t_1+t_2\right)=30\ m/s$

}
\end{ex}

\begin{ex}%book
Từ độ cao $20\ m$, đồng thời thả rơi tự do và ném một vật xuống dưới. Lấy $g=10\ m/s^2$. Khi đó ta phải ném vật vận tốc $v_0$ bằng bao nhiêu để vật này tới mặt đất sớm hơn $1\ s$ so với vật rơi tự do? 
\choice
{\True $15\ m/s$}
{$24\ m/s$}
{$12,5\ m/s$}
{$22,4\ m/s$}
\haidong{20}
\loigiai{
\begin{itemize}
\item Khi vật rơi tự do chạm đất:  $s_1=\dfrac{gt^2}{2}=5t^2=20\Rightarrow t=2\ s$.
\item Vật bị ném chạm đất sớm hơn 1 s nên thời gian vật bị ném chạm đất là $t'=t-1=1\ s$.
\item Khi đó: $s_2=v_0.t'+\dfrac{gt'^2}{2}=h\Rightarrow 20=v_0+5\Rightarrow v_0=15\ m/s$.
\end{itemize}	
}
\end{ex}




\subsubsection*{PHẦN 2. Thí sinh trả lời từ câu 1 đến câu 4. Trong mỗi ý a), b), c), d) ở mỗi câu, thí sinh chọn đúng hoặc sai.}
\setcounter{ex}{0}
\begin{ex}%book
Một quả cầu nhỏ được ném thẳng đứng từ mặt đất lên với vận tốc ban đầu $15\ m/s$. Bỏ qua lực cản không khí và lấy $g=10\ m/s^2$. Chọn gốc tọa độ tại mặt đất, mốc thời gian lúc ném vật, chiều dương có phương thẳng đứng hướng lên trên. 
\choiceTF[t]
{\True Quả cầu sẽ chuyển động chậm dần đều tới độ cao tối đa rồi rơi xuống nhanh dần đều}
{\True Tại thời điểm $t=2 \ s$, quả cầu đang ở vị trí $10\ m$ so với mặt đất}
{Tại thời điểm $t=2 \ s$, quả cầu đang chuyển động theo chiều dương}
{Tại thời điểm $t=2 \ s$, quả cầu đạt độ cao cực đại}
\loigiai{\begin{itemchoice}
\itemch 
\itemch Vị trí của quả cầu sau khi ném 2 s là $x=v_{0} t-\dfrac{1}{2} g t^{2}=15 . 2-\dfrac{1}{2} . 10 . 2^{2}=10\ \ m$
\itemch Vận tốc của quả cầu sau khi ném 2 s là $v=v_{0}-g t=15-10.2=-5\ m / s$. Dấu “-" thể hiện lúc này quả cầu đang chuyển động ngược chiều dương
\itemch Quả cầu đạt độ cao cực đại khi $v^{\prime}=v_{0}-g t=0 \Leftrightarrow 15-10 . t=0 \Rightarrow t=1,5\ s$
\end{itemchoice}
}
\haidong{15}
\end{ex}
\begin{ex}%Tailieuchuan
Một vật được ném thẳng đứng từ mặt đất lên với vận tốc ban đầu $20 \ m/s$. Bỏ qua sức cản không khí. Lấy $g=10 \ m/s^2$.
\choiceTF[t]
{\True Độ cao của vật sau khi ném $1,5\ s$ là $18,75 \ m$} 
{\True Tốc độ của vật sau khi ném $1,5\ s$ là $5 \ m/s$} 
{Độ cao tối đa mà vật có thể đạt được là $30\ m$} 
{\True Thời gian vật chuyển động trong không khí là $4\ s$} 
\loigiai{}
\end{ex}
\begin{ex}%Tailieuchuan
Người ta ném một vật từ mặt đất lên cao theo phương thẳng đứng với vận tốc $4,0\ m/s$. Lấy $g=10 \ m/s^2$. Coi sức cản không khí không đáng kể.
\choiceTF[t]
{Thời gian vật chuyển động khi đạt độ cao cực đại là $0,8\ s$} 
{Thời gian vật chuyển động trong không khí là $0,4 \ s$} 
{\True Độ cao cực đại vật đạt được là $0,8 \ m$} 
{\True Độ cao vật đạt đạt được sau khi ném được $0,2\ s$ là $0,6\ m$} 
\loigiai{}
\end{ex}
\begin{ex}%gpt
Một vật được ném thẳng đứng xuống dưới từ độ cao $20\ m$ với vận tốc ban đầu $v_0 = 5\ m/s$. Lấy $g = 10\ m/s^2$. Bỏ qua sức cản không khí.
\choiceTF[t]
{\True Vận tốc của vật sau $1\ s$ kể từ khi ném là $15\ m/s$}
{Thời gian vật chạm đất là $2\ s$} 
{Vận tốc của vật khi chạm đất là $25\ m/s$}
{Quãng đường vật rơi được trong giây thứ hai là $15\ m$}
\loigiai{
\begin{itemchoice}
\itemch Vận tốc sau $1\ s$: $v = v_0 + gt = 5 + 10 = 15\ m/s$ $\Rightarrow$ mệnh đề 1 đúng.
\itemch Giải phương trình: $h = v_0 t + \dfrac{1}{2}gt^2 \Rightarrow 20 = 5t + 5t^2$ $\Rightarrow$ nghiệm $t = 1{,}57\ s$ $\Rightarrow$ mệnh đề 2 sai.
\itemch Tốc độ chạm đất: $v = \sqrt{v_0^2 + 2gh} = \sqrt{25 + 400} = 25\ m/s$ $\Rightarrow$ mệnh đề 3 đúng.
\itemch Quãng đường giây thứ hai: $s_2 = v_0 \cdot 1 + \dfrac{1}{2}g(2^2 - 1^2) = 5 + \dfrac{1}{2} \cdot 10 \cdot 3 = 5 + 15 = 20\ m$ $\Rightarrow$ mệnh đề 4 sai.
\end{itemchoice}
}
\end{ex}
\subsubsection*{PHẦN 3. Thí sinh trả lời từ câu 1 đến câu 6.}
\setcounter{ex}{0} 
\textit{Sử dụng thông tin sau cho Câu 1 và Câu 2:} Từ điểm $A$ cách mặt đất $4,8\ m$ một vật nhỏ được ném lên cao theo phương thẳng đứng với vận tốc ban đầu $5\ m/s$. Lấy $g=10\ m/s^2$. Bỏ qua sức cản không khí.
\begin{ex}%Tailieuchuan
Độ cao cực đại mà vật đạt được so với mặt đất là bao nhiêu mét?
\shortans[oly]{6,05}
\loigiai{}
\end{ex}
\begin{ex}%Tailieuchuan
 Thời gian vật chuyển động cho đến khi chạm đất là bao nhiêu giây?
\shortans[oly]{1,6}
\end{ex}
\begin{ex}
Một vật được ném thẳng đứng lên cao từ mặt đất với vận tốc ban đầu $v_0 = 20\ m/s$. Bỏ qua sức cản không khí. Lấy $g = 10\ m/s^2$. Độ cao cực đại mà vật đạt được là bao nhiêu mét?
\shortans[oly]{20}
\loigiai{
Ta dùng công thức: $h = \dfrac{v_0^2}{2g} = \dfrac{20^2}{2 \cdot 10} = 20\ m$
}
\end{ex}

\begin{ex}
Một vật rơi tự do từ độ cao $45\ m$ xuống đất. Bỏ qua sức cản không khí. Lấy $g = 10\ m/s^2$. Vật mất bao nhiêu giây để chạm đất?
\shortans[oly]{3}
\loigiai{
Ta dùng công thức: $t = \sqrt{\dfrac{2h}{g}} = \sqrt{\dfrac{2 \cdot 45}{10}} = \sqrt{9} = 3\ s$
}
\end{ex}

\begin{ex}
Một vật được ném thẳng đứng xuống với vận tốc $v_0 = 5\ m/s$ từ độ cao $h = 45\ m$. Lấy $g = 10\ m/s^2$. Tốc độ của vật ngay trước khi chạm đất là bao nhiêu $m/s$ (làm tròn kết quả đến chữ số hàng phần mười)?
\shortans[oly]{30,4}
\loigiai{
Áp dụng công thức: $v = \sqrt{v_0^2 + 2gh} = \sqrt{5^2 + 2 \cdot 10 \cdot 45} = \sqrt{925} \approx 30{,}4\ m/s$
}
\end{ex}

\begin{ex}
Một vật được ném thẳng đứng lên với vận tốc ban đầu $v_0 = 12\ m/s$. Lấy $g = 10\ m/s^2$. Sau bao lâu giây thì vật đạt độ cao cực đại?
\shortans[oly]{1,2}
\loigiai{
Thời gian lên cao cực đại: $t = \dfrac{v_0}{g} = \dfrac{12}{10} = 1{,}2\ s$
}
\end{ex}
\newpage 